\section{The extension and its normalization}\label{sec:geometry}

Throughout the paper, $(X,\omega)$ is a connected compact K\"ahler
surface and $V$ is a holomorphic vector bundle in an exact sequence
\begin{equation}\label{eq:extension}
 0\longrightarrow L\xrightarrow{\iota}V
       \xrightarrow{q}\cI_WM\longrightarrow0.
\end{equation}
Here $L$ and $M$ are line bundles of equal $\omega$-degree,
$W$ is a nonempty zero-dimensional analytic subscheme, and $V$
has rank two and is locally free. Existence of this extension is
part of the hypotheses. Write $\Sigma=\Supp W$ and
$U=X\setminus\Sigma$, and also denote by $q$ the composite
$V\to\cI_WM\hookrightarrow M$.
The positive asymptotic theorem will require $W$ to be reduced
and $L\not\simeq M$. We retain the general setting in this section
and in the residue estimate.

\subsection{Metrics, norms and the fixed inclusion}

Use the determinant identification induced by \eqref{eq:extension}.
Choose Hermitian--Einstein metrics $h_L,h_M$ and put
\begin{equation}\label{eq:fixed-normalization}
 h_{\det}=h_Lh_M,\qquad
 \dd V=\frac{\omega^2}{2},\qquad
 v_0=\int_X\dd V,\qquad
 \kappa=\frac{2\pi\deg_\omega L}{v_0}.
\end{equation}
Degrees are $\int_Xc_1(\,\cdot\,)\wedge\omega$.
Both line curvature contractions equal $\kappa$.
All integral norms use the volume form in
\eqref{eq:fixed-normalization}. On Hermitian endomorphisms the
real inner product is $\tr(st)$, and on skew-Hermitian
endomorphisms it is $-\tr(st)$. Form norms are the ordinary real
K\"ahler norms, extended to complex tensors. In Euclidean complex
coordinates with $\omega_E=(i/2)\sum_j\dd z_j\wedge\dd\bar z_j$,
one consequently has $|\dd z_j|^2=2$.

For a smooth metric $H$ on $V$ with determinant $h_{\det}$, define
\begin{equation}\label{eq:observables}
 \begin{aligned}
 \Phi_H&=i\Lambda_\omega F_H-\kappa I,&
 u_H&=\iota^*H/h_L,\\
 Z_H&=\int_Xu_H\,\dd V,&
 \mathcal E_H&=\int_X
       |D'_{H\otimes h_L^{-1}}\iota|^2\,\dd V .
 \end{aligned}
\end{equation}
The endomorphism $\Phi_H$ is Hermitian and trace-free.
The function $u_H$ is smooth on all of $X$, vanishes exactly on
$\Sigma$, and is positive on $U$. The inclusion in
\eqref{eq:extension} is a fixed holomorphic section of
$V\otimes L^{-1}$; it is not renormalized as time varies.

We use the nonnegative scalar Laplacian $\Delta_0=\dd^*\dd$.
Thus $-\Delta_0=4\sum_j\partial_{z_j}\partial_{\bar z_j}$
in Euclidean coordinates. The metric flow is
\begin{equation}\label{eq:metric-flow}
 H^{-1}\partial_TH=-2\Phi_H,\qquad
 H(0)=H_{\mathrm{in}},\qquad \det H_{\mathrm{in}}=h_{\det}.
\end{equation}
The base, the holomorphic extension and the line metrics remain
fixed along this equation.

\subsection{The qualitative flow input}

We state the standard convergence input in the form used below.
Long-time existence holds for every smooth initial Hermitian
metric on a compact K\"ahler base
\cite[Theorem~6.8.25]{Kobayashi87}.
The determinant solves its scalar heat equation, so the initial
condition in \eqref{eq:metric-flow} preserves $h_{\det}$.
The time in that reference is $t=2T$.

The rank-one factors of \eqref{eq:extension} are stable of equal
slope. Thus $V$ is semistable, with graded sheaf
$G=L\oplus\cI_WM$ and reflexive hull $G^{**}=L\oplus M$.
The corrected convergence theorem of Sibley--Wentworth
\cite[Theorem~1.1 and Remark~1.2, with the appended erratum]{SibleyWentworth15}
identifies the connection limit with
\begin{equation}\label{eq:graded-connection}
                 D_*=D_{h_L}\oplus D_{h_M}.
\end{equation}
More precisely, after identifying the evolving metrics with a
fixed Hermitian bundle, every sequence $T_j\to\infty$ has
subsequences and unitary gauges giving smooth convergence on
compact subsets of $U$ to this connection. The singular set is
$\Sigma$, and its multiplicity at $p$ is
$\dim_\CC\OO_{W,p}$. The corrected multiplicity counts the full
local length, including nonreduced structure. We use this
connection theorem and its point cycle; it supplies no rate on
a shrinking neighbourhood.

The following scalar consequences were established in
\cite[Section~5]{W_Singular}. We include their proof to fix the
normalization and to separate the scalar assertion from a
choice of limiting graded summand.

\begin{proposition}\label{prop:qualitative-flow}
For every solution of \eqref{eq:metric-flow},
\begin{equation}\label{eq:exact-dissipation}
       Z'(T)=-2\mathcal E(T),\qquad
       \sup_Xu(T)\le CZ(T).
\end{equation}
For every compact $K\subset U$,
\begin{equation}\label{eq:normalized-exterior}
          \frac{u(T)}{Z(T)}\longrightarrow\frac1{v_0}
                     \quad\hbox{in }C^\infty(K).
\end{equation}
The constant and convergence times may depend on the initial metric.
\end{proposition}

\begin{proof}
The curvature contraction on $V\otimes L^{-1}$ is $\Phi_H$.
Differentiating the norm of the fixed section and using its
Bochner identity \cite[Proposition~3.1.8]{Kobayashi87} gives
\[
 \partial_Tu=-2\langle\Phi_H\iota,\iota\rangle,\qquad
 -\Delta_0u=2|D'\iota|^2
                         -2\langle\Phi_H\iota,\iota\rangle .
\]
Integration on the compact surface proves the first identity in
\eqref{eq:exact-dissipation}. No singular boundary is introduced:
the section and its finite-time metric are smooth across $\Sigma$.

The mean-curvature heat identity, in the present time convention, is
\begin{equation}\label{eq:mean-curvature-heat}
 (\partial_T+\Delta_0)|\Phi_H|^2=-2|D_H\Phi_H|^2\le0;
\end{equation}
see \cite[Proposition~6.9.1 and equation~(6.9.2)]{Kobayashi87}.
Its maximum principle bounds $\|\Phi_H\|_\infty$ by the initial
value. The scalar Bochner inequality therefore gives
$-\Delta_0u\ge-Cu$. The mean-value estimate on the fixed compact
base yields $\sup u\le C\int u$, as asserted.

Transport the inclusion isometrically to the fixed Hermitian
bundle used in the connection convergence theorem and divide it
by $Z(T)^{1/2}$. The resulting holomorphic section $j_T$ has
integral squared norm one and a uniform pointwise norm bound.
On a subsequence giving smooth connection convergence on $U$,
interior estimates for its holomorphic equation give smooth
subsequential convergence to a holomorphic section of
$(L\oplus M)\otimes L^{-1}$ on $U$. Its coefficients extend
over the isolated points by Hartogs' theorem. The resulting
global bundle has zero contracted curvature, so its integrated
Bochner identity makes the section parallel.

The pointwise bound excludes loss of section mass at the points:
the mass in a union of radius-$r$ coordinate balls is at most
their volume times a fixed constant. First pass to the limit on
the complement and then let $r\downarrow0$. The limiting parallel
section still has squared integral norm one, and hence constant
squared norm $1/v_0$. If the line factors coincide, different
subsequences may give different parallel vectors; their squared
norm is nevertheless the same. Since $|j_T|^2=u/Z$ is unchanged
by the isometries and gauges, a subsequence contradiction proves
the full scalar convergence \eqref{eq:normalized-exterior}.
\end{proof}

\subsection{Local defect data}

In local holomorphic frames compatible with the determinant, the
extension has the presentation
\begin{equation}\label{eq:local-koszul}
             \iota=(f,-g)^{\mathsf t},\qquad q=(g,f).
\end{equation}
The ideal $(f,g)$ has an isolated zero. In the regular local
surface ring it is a regular sequence, so local complete
intersection structure follows from the assumed locally free
extension. Write $m_p=\dim_\CC\OO_{X,p}/(f,g)$.

At a reduced point choose holomorphic coordinates orthonormal
for the base metric at $p$. For any smooth determinant-fixed
reference metric $H_0$ the function $u_0=\iota^*H_0/h_L$ has
quadratic term
\[
                  u_0(z)=z^*Q_pz+O(|z|^3),
                  \qquad Q_p>0 .
\]
The derivative of the inclusion is an isomorphism at a reduced
zero. Its Gram determinant is independent of $H_0$ with fixed
determinant: in the frames of \eqref{eq:local-koszul},
\begin{equation}\label{eq:defect-constants}
 d_p:=\det Q_p
       =|\det\partial(f,g)(p)|^2\frac{h_M(p)}{h_L(p)},
 \qquad
 \lambda_p=\sqrt{d_p},\qquad
 \sigma=\sum_{p\in\Sigma}\lambda_p^{-1}.
\end{equation}
These expressions retain the fixed inclusion, base and line
normalizations. Unit changes of orthonormal coordinates and
compatible frames do not change $d_p$.

For later estimates we need more than a balanced value at the
point: the reference metric must have a prescribed local form.

\begin{lemma}\label{lem:normalized-reference}
If $W$ is reduced, there are K\"ahler normal holomorphic
coordinates at each point, determinant-compatible holomorphic
frames, and one smooth global metric $H_0$ with determinant
$h_{\det}$ such that, on disjoint coordinate balls,
\begin{equation}\label{eq:normal-reference}
 \begin{gathered}
 \iota=\sqrt{\lambda_p}(z_1,-z_2)^{\mathsf t},\qquad
 q=\sqrt{\lambda_p}(z_2,z_1),\\
 H_0=hI,\qquad h=\sqrt{h_Lh_M},\qquad
 u_0=c_p(z)|z|^2,\qquad
 c_p=\lambda_p\sqrt{h_M/h_L}.
 \end{gathered}
\end{equation}
The line weights equal one at the centre and have vanishing
first derivatives. Consequently
$h=1+O(|z|^2)$, $c_p=\lambda_p+O(|z|^2)$ and
$\omega=\omega_E+O(|z|^2)$.
\end{lemma}

\begin{proof}
Choose K\"ahler normal coordinates and normalize the
holomorphic line frames by a constant and the exponential of a
linear function. This fixes their values and first metric jets.
Complete their product to a determinant-compatible frame of $V$.
Let $v(z)$ be the inclusion column and $J=\partial v(0)$.
We first remove its higher terms by a determinant-one frame
change while leaving the base coordinates fixed.

Write $v=A(z)z$, where $A(0)=J$ and $A$ is holomorphic and
invertible on a smaller ball. Put $B=JA^{-1}$ and $d=\det B$.
Then $Bv=Jz$ and $d(0)=1$. For $q_z=(-z_2,z_1)$ choose a
holomorphic column $w$ with $q_zw=d^{-1}-1$; this is possible
since the right side lies in the maximal ideal. The matrix
$C=J(I+wq_z)J^{-1}$ fixes $Jz$ and has determinant $d^{-1}$.
Thus $CB$ has determinant one and sends $v$ to $Jz$.

A unitary change of the base-coordinate phase makes
$\det J=-\lambda_p$. A constant determinant-one fibre change
then makes $J=\sqrt{\lambda_p}\diag(1,-1)$.
Determinant compatibility gives the quotient in
\eqref{eq:normal-reference}. Prescribe $H_0=\sqrt{h_Lh_M}\,I$
on smaller balls. Its determinant is exactly $h_{\det}$ and its
inclusion norm is the one displayed. To extend these prescriptions,
start with any global determinant-fixed metric. Express the local
relative changes by trace-free self-adjoint logarithms, multiply
by cutoffs supported in the disjoint balls, and exponentiate.
The result agrees with the prescribed metrics near each point
and retains the determinant everywhere.
\end{proof}

The hypotheses of this lemma are used only where reducedness is
stated. A nonreduced ideal has a singular derivative and does not
admit this linear local presentation.
